\scanpage{145}
\mspage{12.}
La démonstration de ce théorème sera donnée dans un instant.

\uline{Définition 1}.- \uline{Le topos minimal du théorème} 2 \uline{s'appelle le}
(petit) \uline{topos étale sur} $X$ \uline{et nous dirons qu'un morphisme} $Y/X$
\uline{est étale si} $Y/X$ \uline{appartient au topos étale sur} $X$.

\uline{Remarque 1}.- Il est clair que le morphisme de $X$-préschémas
$S'\longrightarrow S$ de la condition (A') induit un homéomorphisme des espaces
topologiques sous-jacents, de sorte que (A') est bien moins
éxigeant que (A). Dans l'exposé suivant on verra que les $F$ du
topos étale ont en effet des propriétés \inserted{en apparence} plus fortes que (A')
(cf. remarque 3 du th.1 pour la situation analogue pour le topos
de Zariski ).

\uline{Remarque 2}.- La définition ``axiomatique'' du topos étale donnée
ici est évidemment inspirée par la définition ``fonctorielle''
des morphismes étales due à Grothendieck ([SGA 1963, XI,1] et
[EGA IV\qquad] ainsi que [SGA 1960/61,III,3]).Nous renvoyons
le lecteur à [EGA IV\qquad] pour la démonstration des assertions
suivantes:

\uline{Lemme 4}.-\uline{Pour que} $Y\xrightarrow{f}X$ \uline{soit étale, il faut et il suffit que}
\uline{les conditions suivantes soient vérifiées:}

(a) $Y$ \uline{est localement de présentation finie sur} $X$

(b) $Y$ \uline{est plat sur} $X$, \uline{i.e. pour tout} $y\in Y$, $O_y$ \uline{est un} $O_{f(y)}$-
\uline{module plat}.

(c) \uline{Pour chaque} $y\in Y$, $\kappa(y)=O_y/m_y$ \uline{est une algèbre finie}
\uline{séparable sur} $\kappa(x)=O_x/m_x$ (\inserted{$x=f(y)$ et} $m_x$ resp. $m_y$ désigne l'idéal
maximal de l'anneau local $O_x$ resp. $O_y$). \uline{Il revient au même}
\uline{de dire que} $m_y=m_xO_y$ \uline{et que l'extension de corps résiduels} \inserted{$\kappa(y)/\kappa(x)$}
\uline{est une extension finie et séparable}.

\uline{Lemme 5}.- \uline{Pour que} $Y\xrightarrow{f}X$ \uline{soit une immersion ouverte, il faut}
\uline{et il suffit que} $f$ \uline{soit étale et radiciel}. (Rappelons que $f$
radiciel signifie que $f$ est injectif et que pour tout $y\in Y$
l'extension de corps $\kappa(y)/\kappa(f(y))$ est radiciel. On peut
\mspage{144}
